% ==============================================================================
% 05_criteria2_structure.tex — Criterion 2: Programme Structure and Content
% Score: 2
% ==============================================================================

\criterion{2}{โครงสร้างและเนื้อหาหลักสูตร}{Programme Structure and Content}

% ------------------------------------------------------------------------------
\criterionsub{2.1 The specifications of the programme and all its courses are shown to be comprehensive, up-to-date, and made available and communicated to all stakeholders.}

หลักสูตรวิศวกรรมศาสตรมหาบัณฑิต สาขาวิชาวิศวกรรมคอมพิวเตอร์และปัญญาประดิษฐ์ จัดทำข้อกำหนดหลักสูตร (Programme Specification) ในรูปแบบ มคอ.2 ซึ่งได้รับอนุมัติจากสภามหาวิทยาลัยราชภัฏอุตรดิตถ์ ครั้งที่ 4/2568 วันที่ 4 เมษายน พ.ศ.~2568 ก่อนเปิดรับนักศึกษาในภาคการศึกษาที่ 2/2568

มคอ.2 ฉบับที่ผ่านสภามีเนื้อหาครอบคลุมตามองค์ประกอบที่กำหนด ได้แก่ ชื่อหลักสูตรและปริญญา หน่วยงานที่รับผิดชอบ ปรัชญาและวัตถุประสงค์ ผลการเรียนรู้ที่คาดหวัง (PLOs) โครงสร้างรายวิชา แผนการศึกษา เกณฑ์การรับนักศึกษา คุณสมบัติบัณฑิต และข้อมูลอาจารย์ประจำหลักสูตร ข้อกำหนดดังกล่าวได้รับการปรับปรุงล่าสุดก่อนเปิดรับนักศึกษา โดยมีการแก้ไขตามข้อเสนอแนะของสภามหาวิทยาลัย เช่น การปรับรหัสวิชากลุ่ม AI/ML จาก prefix 701 เป็น 409 ตามมติสภา (อ้างอิง~\hyperlink{ref:AUNQA-2-1}{AUNQA-2-1})

ข้อกำหนดรายวิชา (Course Specification) ในรูปแบบ มคอ.3 ได้จัดทำครบทุกรายวิชาที่เปิดสอนในปีการศึกษา 2568 โดยมคอ.3 แต่ละฉบับระบุ: ชื่อและรหัสวิชา วิชาที่ต้องเรียนมาก่อน ผลการเรียนรู้รายวิชา (CLOs) กิจกรรมการเรียนการสอน วิธีการวัดและประเมินผล ตารางการสอน และวันที่จัดทำ ข้อมูลมคอ.3 ทุกวิชาจัดเก็บอย่างเป็นระบบและพร้อมให้ตรวจสอบ (อ้างอิง~\hyperlink{ref:AUNQA-2-4}{AUNQA-2-4})

หลักสูตรเผยแพร่ข้อมูลหลักสูตรผ่านช่องทางหลายรูปแบบ ได้แก่ เว็บไซต์มหาวิทยาลัยราชภัฏอุตรดิตถ์ เว็บไซต์คณะเทคโนโลยีอุตสาหกรรม Facebook Page ``วิศวกรรมคอมพิวเตอร์และปัญญาประดิษฐ์'' (491 followers ณ พฤษภาคม 2569) เอกสารแนะนำหลักสูตรสำหรับผู้สมัคร และผ่านกิจกรรมปฐมนิเทศนักศึกษา ซึ่งจัดขึ้นเมื่อวันที่ 12 ตุลาคม พ.ศ.~2568 โดยมีรองอธิการบดีและคณบดีคณะเทคโนโลยีอุตสาหกรรมเข้าร่วม นอกจากนี้หลักสูตรได้ประชาสัมพันธ์ที่วิทยาลัยเทคนิคอุตรดิตถ์ (2 ตุลาคม 2568) เพื่อเข้าถึงกลุ่มเป้าหมายที่หลากหลาย

% ------------------------------------------------------------------------------
\criterionsub{2.2 The design of the curriculum is shown to be constructively aligned with achieving the expected learning outcomes.}

หลักสูตรออกแบบโดยยึดหลัก Outcome-Based Education (OBE) และ Backward Curriculum Design (BCD) ตามแนวทางของ Biggs (2014) กล่าวคือ การกำหนดรายวิชาและลำดับการเรียนรู้เริ่มจาก PLOs ที่ต้องการให้นักศึกษาบรรลุเมื่อสำเร็จการศึกษา แล้วจึงออกแบบ CLOs กิจกรรมการเรียนการสอน (Teaching and Learning Activities: TLAs) และการประเมินผลให้สอดคล้องเชื่อมโยงกัน (Constructive Alignment)

ในระดับรายวิชา มคอ.3 ทุกฉบับกำหนด CLOs ที่ map โดยตรงกับ PLO ระบุกิจกรรมการเรียนการสอนที่รองรับ CLO แต่ละข้อ และกำหนดวิธีการประเมินผลที่สอดคล้องกัน โดยใช้รูบริกสี่ระดับคุณภาพสำหรับการประเมิน CLO ทุกข้อ และโครงสร้างคะแนน 30/40/30 (ก่อนกลางภาค/กลางภาค/ปลายภาค) สำหรับทุกรายวิชาบังคับ

ตาราง~\ref{tbl:clo_plo_core} แสดงการ mapping ระหว่าง CLO กับ PLO ของวิชาบังคับหลักทั้ง 4 วิชา ซึ่งครอบคลุม PLO ทุกข้อ โดย PLO5 ครอบคลุมผ่านวิชา 7015906 ระเบียบวิธีวิจัย (CLO5) เท่านั้นในกลุ่มวิชาบังคับ ซึ่งนับเป็นหลักฐานของการออกแบบเชิงกลยุทธ์ที่มอบหมายความรับผิดชอบแก่ PLO5 ไว้ที่วิชาที่เหมาะสมที่สุด

\begin{table}[H]
  \centering
  \caption{CLO--PLO Mapping วิชาเฉพาะด้านบังคับ 4 วิชาหลัก}
  \label{tbl:clo_plo_core}
  \begingroup
  \small
  \setlength{\tabcolsep}{4pt}
  \renewcommand{\arraystretch}{1.2}
  \begin{tabularx}{\linewidth}{|>{\raggedright\arraybackslash}X|c|c|c|c|c|}
    \hline
    \textbf{รายวิชา} & \textbf{PLO1} & \textbf{PLO2} & \textbf{PLO3} & \textbf{PLO4} & \textbf{PLO5} \\
    \hline
    4095101 ปัญญาประดิษฐ์และการเรียนรู้ของเครื่อง & CLO1,2,5 & & CLO3,4 & & \\
    \hline
    4095102 การเขียนโปรแกรมขั้นสูงฯ & CLO1,2,5 & CLO2,3 & CLO3,5 & CLO4 & \\
    \hline
    7015101 เทคโนโลยีอุบัติใหม่ฯ & CLO1,2,4,5 & CLO3 & & & \\
    \hline
    7015906 ระเบียบวิธีวิจัยฯ & & CLO1 & CLO2,3 & CLO4 & CLO5 \\
    \hline
    \textbf{รวม PLO ที่ครอบคลุม} & \checkmark & \checkmark & \checkmark & \checkmark & \checkmark \\
    \hline
  \end{tabularx}
  \endgroup
\end{table}

% ------------------------------------------------------------------------------
\criterionsub{2.3 The design of the curriculum is shown to include feedback from stakeholders, especially external stakeholders.}

หลักสูตรดำเนินการจัดเก็บข้อมูลจากผู้มีส่วนได้ส่วนเสียเพื่อนำมาประกอบการออกแบบหลักสูตรก่อนการยกร่าง โดยมีการระบุกลุ่มผู้มีส่วนได้ส่วนเสียและวิธีการรับข้อมูลที่หลากหลาย

ข้อมูลที่ได้จากผู้มีส่วนได้ส่วนเสียสะท้อนให้เห็นในการออกแบบหลักสูตรในหลายมิติ ได้แก่ (1) ด้านเนื้อหา: มีวิชา IoT Smart Farming, AI สำหรับงาน Healthcare และ Drone AI ซึ่งตอบสนองความต้องการขององค์กรในท้องถิ่น (2) ด้านรูปแบบการเรียน: เปิดสอนในวันเสาร์--อาทิตย์ (ร้อยละ~83.9 ของกลุ่มเป้าหมายต้องการ) (3) ด้านทักษะ: เน้นทักษะการเขียนโปรแกรม MLOps และ Git/GitHub ตามความต้องการของสถานประกอบการ

ทั้งนี้ ในฐานะหลักสูตรใหม่ที่เปิดรับนักศึกษารุ่นแรกในภาค 2/2568 หลักสูตรได้วางระบบเพื่อรับ Feedback จากผู้มีส่วนได้ส่วนเสียที่สัมผัสกับหลักสูตรจริงอย่างต่อเนื่อง โดยจะดำเนินการสำรวจประสบการณ์นักศึกษาและผู้ใช้บัณฑิตอย่างเป็นทางการในรอบการทบทวนประจำปี

% ------------------------------------------------------------------------------
\criterionsub{2.4 The contribution made by each course in achieving the expected learning outcomes is shown to be clear.}

หลักสูตรจัดทำตาราง Curriculum Mapping ที่แสดงให้เห็นอย่างชัดเจนว่ารายวิชาใดรับผิดชอบหลักต่อ PLO ใด (primary contribution, $\bullet$) และรายวิชาใดสนับสนุนเสริม PLO นั้น (secondary contribution, $\circ$)

\begin{table}[H]
  \centering
  \caption{Curriculum Mapping — ความรับผิดชอบของรายวิชาต่อ PLO}
  \label{tbl:curriculum_mapping}
  \begingroup
  \footnotesize
  \setlength{\tabcolsep}{3pt}
  \renewcommand{\arraystretch}{1.15}
  \begin{tabularx}{\linewidth}{|>{\raggedright\arraybackslash}X|c|c|c|c|c|}
    \hline
    \multicolumn{1}{|c|}{\textbf{รายวิชา}} & \textbf{PLO1} & \textbf{PLO2} & \textbf{PLO3} & \textbf{PLO4} & \textbf{PLO5} \\
    \hline
    \multicolumn{6}{|l|}{\textit{วิชาเฉพาะด้านบังคับ}} \\
    \hline
    4095101 ปัญญาประดิษฐ์และการเรียนรู้ของเครื่อง & $\bullet$ & & $\bullet$ & & \\
    \hline
    4095102 การเขียนโปรแกรมขั้นสูงสำหรับ ML & $\bullet$ & $\circ$ & $\circ$ & $\bullet$ & \\
    \hline
    7015101 เทคโนโลยีอุบัติใหม่ฯ & $\bullet$ & $\circ$ & & & \\
    \hline
    7015906 ระเบียบวิธีวิจัยฯ & & $\bullet$ & $\bullet$ & $\circ$ & $\bullet$ \\
    \hline
    7015907 สัมมนาวิศวกรรมคอมพิวเตอร์และ AI & & $\bullet$ & $\circ$ & $\circ$ & $\bullet$ \\
    \hline
    7015908 หัวข้อพิเศษฯ & $\bullet$ & $\circ$ & & $\circ$ & \\
    \hline
    \multicolumn{6}{|l|}{\textit{วิทยานิพนธ์ (แผน 1)}} \\
    \hline
    7015901--7015903 วิทยานิพนธ์ 1--3 & $\bullet$ & $\bullet$ & $\bullet$ & $\bullet$ & $\bullet$ \\
    \hline
    \multicolumn{6}{|l|}{\textit{สารนิพนธ์ (แผน 2)}} \\
    \hline
    7015904--7015905 สารนิพนธ์ 1--2 & $\bullet$ & $\bullet$ & $\bullet$ & $\bullet$ & $\circ$ \\
    \hline
    \multicolumn{6}{|l|}{\textit{วิชาเสริม (ไม่นับหน่วยกิต)}} \\
    \hline
    1555101 ภาษาอังกฤษสำหรับบัณฑิตศึกษา & & $\circ$ & & & $\bullet$ \\
    \hline
    1506001 การเขียนวิทยานิพนธ์หรือสารนิพนธ์ & & $\circ$ & $\circ$ & $\circ$ & $\bullet$ \\
    \hline
  \end{tabularx}
  \endgroup
  \begin{flushleft}
    \footnotesize $\bullet$ = ผลักดันหลัก (primary); $\circ$ = ผลักดันรอง (secondary)
  \end{flushleft}
\end{table}

วิชาเลือก 22 วิชา (รวม 6--12 หน่วยกิต ขึ้นอยู่กับแผนการศึกษา) มีการ map กลับ PLO1--PLO3 เป็นหลัก เพื่อขยายและเจาะลึกองค์ความรู้เฉพาะทางตามความสนใจของนักศึกษา ข้อมูล CLO--PLO mapping ของรายวิชาเลือกทุกวิชาปรากฏใน เอกสาร \hyperlink{ref:AUNQA-2-2}{AUNQA-2-2} (รวมเล่ม CLO รายวิชา)

% ------------------------------------------------------------------------------
\criterionsub{2.5 The curriculum to show that all its courses are logically structured, properly sequenced, and are integrated.}

หลักสูตรออกแบบโครงสร้างรายวิชาให้มีลำดับการเรียนรู้ที่เป็นระบบตามหลัก Learning Progression จาก Basic ไปสู่ Intermediate และ Specialized ดังนี้

\textbf{ชั้นปีที่ 1 ภาคการศึกษาที่ 1} (พื้นฐาน): วิชา 4095101 ปัญญาประดิษฐ์และการเรียนรู้ของเครื่อง และ 7015906 ระเบียบวิธีวิจัยทางวิทยาศาสตร์และวิศวกรรมศาสตร์ ถูกกำหนดให้เรียนก่อนโดยเป็น Foundation ทั้งด้านทักษะ AI/ML และทักษะการวิจัย

\textbf{ชั้นปีที่ 1 ภาคการศึกษาที่ 2} (ประยุกต์): วิชา 4095102 การเขียนโปรแกรมขั้นสูงสำหรับ ML ต่อยอดทักษะการเขียนโปรแกรมสำหรับ ML ในสภาพแวดล้อมจริง และ 7015101 เทคโนโลยีอุบัติใหม่ฯ เพิ่มมุมมองเทคโนโลยีสมัยใหม่ รวมถึง 7015907 สัมมนาวิศวกรรมฯ ที่บูรณาการทักษะวิจัยและการสื่อสาร

\textbf{ชั้นปีที่ 2} (เชี่ยวชาญและบูรณาการ): วิชาเลือกและวิทยานิพนธ์/สารนิพนธ์เป็นพื้นที่ Specialization โดย 7015908 หัวข้อพิเศษฯ ช่วยเจาะลึกตามความสนใจ ก่อนที่นักศึกษาจะบูรณาการความรู้ทั้งหมดในวิทยานิพนธ์หรือสารนิพนธ์

\textbf{ตัวอย่าง Integration — การถ่ายทอดองค์ความรู้สู่ชุมชน}: วิชา 7015101 เทคโนโลยีอุบัติใหม่ฯ กำหนดให้สัปดาห์ที่~10 เป็น IEL Community Workshop ที่นักศึกษาบูรณาการความรู้ข้ามวิชา ได้แก่ หลักการ AI/ML จากวิชา 4095101, ทักษะการเขียนโปรแกรมจากวิชา 4095102 และกรอบเทคโนโลยีอุบัติใหม่จากวิชา 7015101 เพื่อนำไปถ่ายทอดให้กับกลุ่มเป้าหมายภายนอกในฐานะ \textbf{วิทยากร} โดยมีกิจกรรมที่เกิดขึ้นจริงในปีการศึกษา~2568--2569 ดังนี้

\begin{itemize}
  \item \textbf{อบรม AI สำหรับครูประจำการ} (26 มีนาคม 2569): นักศึกษาระดับบัณฑิตศึกษาเป็นวิทยากรหลักในการอบรม ``การใช้งาน AI เพื่อการศึกษา'' ให้กับครูประจำการจากโรงเรียนเทศบาล~3 แห่ง ได้แก่ โรงเรียนเทศบาลท่าอิฐ โรงเรียนวัดหนองผา และโรงเรียนวัดเกษมจิตตาราม สะท้อนการนำความรู้จากชั้นเรียนสู่บริบทจริงของผู้สอนระดับประถม-มัธยมศึกษา

  \item \textbf{อบรม AI สำหรับสหกรณ์ออมทรัพย์ครูอุตรดิตถ์} (4 และ 18 ตุลาคม 2568): นักศึกษาร่วมกับอาจารย์ผู้รับผิดชอบหลักสูตรถ่ายทอดการใช้งาน AI และ IoT Smart Farming~2 ครั้ง ให้กับบุคลากรสหกรณ์ออมทรัพย์ครูอุตรดิตถ์ ซึ่งส่วนใหญ่เป็นครูประจำการในพื้นที่ สะท้อนการบูรณาการ PLO1 (ซอฟต์แวร์/ฮาร์ดแวร์) และ PLO4 (จรรยาบรรณ/ความรับผิดชอบต่อสังคม) ในบริบทการเผยแพร่ความรู้สู่ชุมชน
\end{itemize}

กิจกรรมทั้งสองแสดงให้เห็นว่าหลักสูตรสามารถ Integration ได้ในระดับที่นักศึกษาไม่เพียงแต่ใช้ความรู้เชิงทฤษฎี แต่สามารถสื่อสาร อธิบาย และถ่ายทอดองค์ความรู้ AI ให้กับผู้ที่ไม่มีพื้นฐานด้านคอมพิวเตอร์ได้จริง ซึ่งตรงกับ PLO5 (ทักษะการสื่อสารและการเรียนรู้ตลอดชีวิต) และ PLO2 (การพัฒนาระบบสำหรับท้องถิ่น)


% ------------------------------------------------------------------------------
\criterionsub{2.6 The curriculum to have option(s) for students to pursue major and/or minor specialisations.}

หลักสูตรจัดให้นักศึกษามีทางเลือกในการเรียนรู้ 2 ระดับ ได้แก่

\textbf{ทางเลือกระดับแผนการศึกษา}: นักศึกษาเลือกระหว่าง แผน 1 (วิชาการ) ที่เน้นการทำวิทยานิพนธ์เต็มรูปแบบ 12 หน่วยกิต เหมาะสำหรับผู้ต้องการมุ่งสู่การทำวิจัยระดับสูงหรือศึกษาต่อปริญญาเอก และ แผน 2 (วิชาชีพ) ที่เน้นการทำสารนิพนธ์ 6 หน่วยกิต ร่วมกับวิชาเลือกเพิ่มเติม 6 หน่วยกิต เหมาะสำหรับผู้ต้องการนำความรู้ไปประยุกต์ใช้ในองค์กรหรืออุตสาหกรรม

\textbf{ทางเลือกระดับวิชาเลือก}: วิชาเลือก 22 วิชา ครอบคลุมกลุ่มความเชี่ยวชาญหลายด้าน เช่น Computer Vision, Natural Language Processing, IoT and Embedded Systems, Data Engineering, Cybersecurity, Cloud Computing และ Bioinformatics นักศึกษาสามารถเลือกเรียนได้ 2--4 วิชาตามแผน โดยสาระของวิชาเลือกสอดคล้องกับทิศทางวิทยานิพนธ์และความต้องการของตลาดแรงงาน

% ------------------------------------------------------------------------------
\criterionsub{2.7 The programme to show that its curriculum is reviewed periodically following an established procedure and that it remains up-to-date and relevant to industry.}

% อัปเดต 2026-05-22: เพิ่มหลักฐาน Benchmarking หลักสูตรอ้างอิง (จุฬาฯ มจพ. มช.) จากประชุม 2/2566

หลักสูตรวางระบบการทบทวนหลักสูตรอย่างเป็นระบบตาม PDCA cycle โดยอ้างอิงกรอบระยะเวลาที่กำหนดโดย สป.อว. (ปรับปรุงเล็กน้อยรายปี/ปรับปรุงใหญ่ทุก 5 ปี)

\textbf{Curriculum Benchmarking (ระยะพัฒนาหลักสูตร)}: ในการประชุมคณะกรรมการบริหารหลักสูตรครั้งที่~2/2566 (20~ตุลาคม~2566) หลักสูตรได้ศึกษาและเปรียบเทียบโครงสร้างหลักสูตรอ้างอิงจากมหาวิทยาลัยชั้นนำในประเทศ ได้แก่ \textbf{จุฬาลงกรณ์มหาวิทยาลัย, มหาวิทยาลัยเทคโนโลยีพระจอมเกล้าพระนครเหนือ (มจพ.) และมหาวิทยาลัยเชียงใหม่ (มช.)} ผลการเปรียบเทียบนำไปสู่การออกแบบโครงสร้างรายวิชาบังคับ~6 วิชา (~18 หน่วยกิต) และวิชาเลือก~22 วิชา ที่สอดคล้องกับมาตรฐานหลักสูตรระดับสากล (อ้างอิง~\hyperlink{ref:AUNQA-2-8}{AUNQA-2-8})

\textbf{กระบวนการทบทวนหลักสูตร (PDCA)}: ในฐานะหลักสูตรที่เพิ่งเปิดดำเนินการในปีการศึกษา~2568 กระบวนการทบทวนหลักสูตรอย่างเป็นทางการรอบแรกจะดำเนินการเมื่อสิ้นปีการศึกษา~2568 อย่างไรก็ตาม หลักสูตรได้วางโครงสร้างของระบบทบทวนไว้ดังนี้

\begin{enumerate}
  \item \textbf{การทบทวนรายวิชาประจำปี} ผ่านการจัดทำ มคอ.5 เมื่อสิ้นสุดแต่ละภาคการศึกษา อาจารย์ผู้สอนรายงานผลการประเมิน วิเคราะห์ปัญหา และเสนอแนวทางปรับปรุง (อ้างอิง~\hyperlink{ref:AUNQA-2-7}{AUNQA-2-7}) หลักฐานรอบแรก: มคอ.5 ทั้ง 4 วิชา ลงนาม 30~เมษายน~2569 ระบุ corrective actions รอบถัดไป
  \item \textbf{การทบทวนหลักสูตรประจำปี} โดยคณะกรรมการบริหารหลักสูตร พิจารณาผลการดำเนินการของรายวิชาทั้งหมด ข้อมูลการจ้างงาน (เมื่อมีบัณฑิต) และแนวโน้มเทคโนโลยี เพื่อพิจารณาปรับปรุงรายละเอียดรายวิชาหรือกิจกรรมการเรียนการสอน
  \item \textbf{การปรับปรุงหลักสูตรทุก 5 ปี} ตามกำหนดของ สป.อว. โดยผ่านกระบวนการรับฟังผู้มีส่วนได้ส่วนเสีย วิเคราะห์ตลาดแรงงาน และผ่านกลไกอนุมัติของมหาวิทยาลัย
\end{enumerate}

\textbf{ความทันสมัยของเนื้อหา (Currency)}: ในปีการศึกษา~2568 หลักสูตรตรวจสอบความทันสมัยของเนื้อหาโดยการบูรณาการหัวข้อ AI สมัยใหม่ เช่น Generative AI, Large Language Models และ AI Ethics ไว้ใน syllabus ของวิชาบังคับ และจัดอบรมระยะสั้น (GenAI Short Courses~4 คอร์ส เมษายน--พฤษภาคม~2569) เพื่อเสริมทักษะที่ตลาดต้องการ ซึ่งแสดงให้เห็นถึงความสามารถในการปรับตัวตามการเปลี่ยนแปลงของเทคโนโลยีอย่างรวดเร็ว นอกจากนี้นักศึกษา (n=4) ยังเสนอให้เพิ่มเนื้อหา Big Data/Cloud Computing (4/4 คน) และ AI Security (3/4 คน) ซึ่งหลักสูตรจะนำไปพิจารณาในรอบ PDCA ปีการศึกษา~2569

\begin{strengthbox}
หลักสูตรออกแบบตามหลัก OBE + Backward Curriculum Design โดยมี มคอ.2 ฉบับผ่านสภา มคอ.3 ทุกรายวิชา ตาราง Curriculum Mapping พร้อม primary/secondary designation ครบถ้วน โครงสร้างการเรียนรู้เป็นระบบจาก Basic ถึง Specialized พร้อมตัวอย่าง Integration (IEL community workshop) และมีทางเลือก 2 แผนการศึกษา + วิชาเลือก 22 วิชา
\end{strengthbox}

\begin{weaknessbox}
R2.3 Feedback จากนักศึกษาที่สัมผัสหลักสูตรจริง: ผลสำรวจความพึงพอใจ $n=4$ (ปีการศึกษา~2568) ให้คะแนนโครงสร้าง/เนื้อหา~4.75/5.00 แต่นี่เป็นปีแรก ยังต้องรวบรวม feedback เชิงลึกจากนักศึกษาเมื่อสิ้นปีการศึกษาผ่าน มคอ.5; R2.7 การทบทวนรายปีรอบแรกยังไม่แล้วเสร็จ; CLO--PLO mapping ของวิชาเลือก~22 วิชาอยู่ใน \hyperlink{ref:AUNQA-2-2}{AUNQA-2-2} ยังไม่ได้แสดงในตารางหลัก
\end{weaknessbox}

\begin{selfassessmentbox}{2}{หลักสูตรมีโครงสร้างและเนื้อหาที่ออกแบบอย่างเป็นระบบตาม OBE/BCD มี Curriculum Mapping คุณภาพสูงพร้อม primary/secondary designation ทางเลือกครบถ้วน และมีหลักฐาน currency ของเนื้อหา (GenAI Short Courses) ผลสำรวจความพึงพอใจนักศึกษา $n=4$ ปีการศึกษา~2568 ให้คะแนนโครงสร้าง~4.75/5.00 สอดคล้องกับการออกแบบที่ดี แต่ยังต้องการ Feedback เชิงลึกผ่าน มคอ.5 และผลการทบทวนรายปีรอบแรก}
\end{selfassessmentbox}

\begin{evidencelist}
\evidence{AUNQA-2-1}{เล่ม มคอ.2 หลักสูตรวิศวกรรมคอมพิวเตอร์และปัญญาประดิษฐ์ ฉบับผ่านสภา\ \href{https://syweb.kidcbc.work/Eviden_aunqa68/\%E0\%B8\%A1\%E0\%B8\%84\%E0\%B8\%AD\%202\%20\%E0\%B8\%AB\%E0\%B8\%A5\%E0\%B8\%B1\%E0\%B8\%81\%E0\%B8\%AA\%E0\%B8\%B9\%E0\%B8\%95\%E0\%B8\%A3\%E0\%B8\%A7\%E0\%B8\%B4\%E0\%B8\%A8\%E0\%B8\%A7\%E0\%B8\%81\%E0\%B8\%A3\%E0\%B8\%A3\%E0\%B8\%A1\%E0\%B8\%84\%E0\%B8\%AD\%E0\%B8\%A1\%E0\%B8\%9E\%E0\%B8\%B4\%E0\%B8\%A7\%E0\%B9\%80\%E0\%B8\%95\%E0\%B8\%AD\%E0\%B8\%A3\%E0\%B9\%8C\%E0\%B9\%81\%E0\%B8\%A5\%E0\%B8\%B0\%E0\%B8\%9B\%E0\%B8\%B1\%E0\%B8\%8D\%E0\%B8\%8D\%E0\%B8\%B2\%E0\%B8\%9B\%E0\%B8\%A3\%E0\%B8\%B0\%E0\%B8\%94\%E0\%B8\%B4\%E0\%B8\%A9\%E0\%B8\%90\%E0\%B9\%8C\%20\%E0\%B8\%9C\%E0\%B9\%88\%E0\%B8\%B2\%E0\%B8\%99\%E0\%B8\%AA\%E0\%B8\%A0\%E0\%B8\%B2\%202\%20\%E0\%B9\%80\%E0\%B8\%9E\%E0\%B8\%B4\%E0\%B9\%88\%E0\%B8\%A1\%20Comment\%20\%E0\%B8\%AA\%E0\%B8\%A0\%E0\%B8\%B2\%20\%E0\%B9\%81\%E0\%B8\%A5\%E0\%B8\%B0\%20\%E0\%B8\%A1\%E0\%B8\%95\%E0\%B8\%B4\%E0\%B8\%AA\%E0\%B8\%A0\%E0\%B8\%B2.pdf}{\texttt{[PDF]}}}
\evidence{AUNQA-2-2}{รวมเล่ม CLO รายวิชา (ทุก 33 วิชา) พร้อม CLO--PLO Mapping รายวิชา\ \href{https://syweb.kidcbc.work/Eviden_aunqa68/\%E0\%B8\%A3\%E0\%B8\%A7\%E0\%B8\%A1\%E0\%B9\%80\%E0\%B8\%A5\%E0\%B9\%88\%E0\%B8\%A1\%20CLO\%20\%E0\%B8\%A3\%E0\%B8\%B2\%E0\%B8\%A2\%E0\%B8\%A7\%E0\%B8\%B4\%E0\%B8\%8A\%E0\%B8\%B2.pdf}{\texttt{[PDF]}}}
\evidence{AUNQA-2-3}{Facebook Page วิศวกรรมคอมพิวเตอร์และปัญญาประดิษฐ์ (495 followers) --- เผยแพร่ข้อมูลหลักสูตร}
\evidence{AUNQA-2-4}{มคอ.3 รายวิชาบังคับทุกวิชา (4095101, 4095102, 7015101, 7015906, 7015907, 7015908)\ \href{https://syweb.kidcbc.work/Eviden_aunqa68/\%E0\%B8\%A1\%E0\%B8\%84\%E0\%B8\%AD3/AUN3_4095101.pdf}{\texttt{[PDF]}}}
\evidence{AUNQA-2-7-1}{รายงานการประชุมคณะกรรมการประจำคณะเทคโนโลยีอุตสาหกรรม (เสนอ/เห็นชอบหลักสูตร) ปีการศึกษา~2568\ \href{https://syweb.kidcbc.work/Eviden_aunqa68/\%E0\%B8\%AB\%E0\%B8\%A5\%E0\%B8\%B1\%E0\%B8\%81\%E0\%B8\%90\%E0\%B8\%B2\%E0\%B8\%99\%E0\%B8\%81\%E0\%B8\%B2\%E0\%B8\%A3\%E0\%B8\%9B\%E0\%B8\%A3\%E0\%B8\%B0\%E0\%B8\%8A\%E0\%B8\%B8\%E0\%B8\%A1/02_\%E0\%B8\%A3\%E0\%B8\%B0\%E0\%B8\%94\%E0\%B8\%B1\%E0\%B8\%9A\%E0\%B8\%84\%E0\%B8\%93\%E0\%B8\%B0/AUNQA-2-7-1_\%E0\%B8\%9B\%E0\%B8\%A3\%E0\%B8\%B0\%E0\%B8\%8A\%E0\%B8\%B8\%E0\%B8\%A1\%E0\%B8\%81\%E0\%B8\%A3\%E0\%B8\%A3\%E0\%B8\%A1\%E0\%B8\%81\%E0\%B8\%B2\%E0\%B8\%A3\%E0\%B8\%9B\%E0\%B8\%A3\%E0\%B8\%B0\%E0\%B8\%88\%E0\%B8\%B3\%E0\%B8\%84\%E0\%B8\%93\%E0\%B8\%B0_\%E0\%B9\%80\%E0\%B8\%AA\%E0\%B8\%99\%E0\%B8\%AD\%E0\%B8\%AB\%E0\%B8\%A5\%E0\%B8\%B1\%E0\%B8\%81\%E0\%B8\%AA\%E0\%B8\%B9\%E0\%B8\%95\%E0\%B8\%A3_2568.pdf}{\texttt{[PDF]}}}
\evidence{AUNQA-2-7}{มคอ.5 รายวิชาบังคับ 4 วิชา ลงนาม 30 เม.ย. 2569 (7015101, 7015906 ภาค 2/2568; 4095101, 4095102 ภาค 1/2568) --- ข้อมูลหลักฐานที่ต้องการ/มคอ5\_รายวิชาบังคับ/AUN3\_*\_มคอ5.pdf\ \href{https://syweb.kidcbc.work/Eviden_aunqa68/\%E0\%B8\%A1\%E0\%B8\%84\%E0\%B8\%AD5/AUN3_4095101.pdf}{\texttt{[PDF]}}}
\evidence{AUNQA-2-8}{รายงานการประชุม 2/2566 (20 ต.ค. 2566) — การศึกษาหลักสูตรอ้างอิง: จุฬาลงกรณ์มหาวิทยาลัย, มจพ., มช. เพื่อ Benchmark โครงสร้างหลักสูตร วศ.ม. CPE\&AI\ \href{https://syweb.kidcbc.work/Eviden_aunqa68/\%E0\%B8\%AB\%E0\%B8\%A5\%E0\%B8\%B1\%E0\%B8\%81\%E0\%B8\%90\%E0\%B8\%B2\%E0\%B8\%99\%E0\%B8\%81\%E0\%B8\%B2\%E0\%B8\%A3\%E0\%B8\%9B\%E0\%B8\%A3\%E0\%B8\%B0\%E0\%B8\%8A\%E0\%B8\%B8\%E0\%B8\%A1/01_\%E0\%B8\%A3\%E0\%B8\%B0\%E0\%B8\%94\%E0\%B8\%B1\%E0\%B8\%9A\%E0\%B8\%AB\%E0\%B8\%A5\%E0\%B8\%B1\%E0\%B8\%81\%E0\%B8\%AA\%E0\%B8\%B9\%E0\%B8\%95\%E0\%B8\%A3/\%E0\%B8\%81\%E0\%B8\%A1\%E0\%B8\%A8_2566-10-20_\%E0\%B8\%A2\%E0\%B8\%81\%E0\%B8\%A3\%E0\%B9\%88\%E0\%B8\%B2\%E0\%B8\%87\%E0\%B8\%AB\%E0\%B8\%A5\%E0\%B8\%B1\%E0\%B8\%81\%E0\%B8\%AA\%E0\%B8\%B9\%E0\%B8\%95\%E0\%B8\%A3_\%E0\%B8\%84\%E0\%B8\%A3\%E0\%B8\%B1\%E0\%B9\%89\%E0\%B8\%87\%E0\%B8\%97\%E0\%B8\%B5\%E0\%B9\%882.pdf}{\texttt{[PDF]}}}
\evidence{AUNQA-C6-STSAT-2568}{แบบสำรวจความพึงพอใจนักศึกษาและผู้มีส่วนได้ส่วนเสีย CPE\&AI ปีการศึกษา~2568 ($n=4$): คะแนนความพึงพอใจโครงสร้าง/เนื้อหาเฉลี่ย~4.75/5.00 --- student\_satisfaction\_survey\_2568.csv \href{https://docs.google.com/forms/d/1xpfb6AfYkbOkidJNh6ci95ux4ZLYzjvpUuqARYjhr5k/viewanalytics}{\texttt{[Google Report]}}}
\end{evidencelist}
